% Options for packages loaded elsewhere
\PassOptionsToPackage{unicode}{hyperref}
\PassOptionsToPackage{hyphens}{url}
\PassOptionsToPackage{dvipsnames,svgnames*,x11names*}{xcolor}
%
\documentclass[
  a4paper,
  dvipsnames]{tufte-book}
\usepackage{lmodern}
\usepackage{amssymb,amsmath}
\usepackage{ifxetex,ifluatex}
\ifnum 0\ifxetex 1\fi\ifluatex 1\fi=0 % if pdftex
  \usepackage[T1]{fontenc}
  \usepackage[utf8]{inputenc}
  \usepackage{textcomp} % provide euro and other symbols
\else % if luatex or xetex
  \usepackage{unicode-math}
  \defaultfontfeatures{Scale=MatchLowercase}
  \defaultfontfeatures[\rmfamily]{Ligatures=TeX,Scale=1}
\fi
% Use upquote if available, for straight quotes in verbatim environments
\IfFileExists{upquote.sty}{\usepackage{upquote}}{}
\IfFileExists{microtype.sty}{% use microtype if available
  \usepackage[]{microtype}
  \UseMicrotypeSet[protrusion]{basicmath} % disable protrusion for tt fonts
}{}
\makeatletter
\@ifundefined{KOMAClassName}{% if non-KOMA class
  \IfFileExists{parskip.sty}{%
    \usepackage{parskip}
  }{% else
    \setlength{\parindent}{0pt}
    \setlength{\parskip}{6pt plus 2pt minus 1pt}}
}{% if KOMA class
  \KOMAoptions{parskip=half}}
\makeatother
\usepackage{xcolor}
\IfFileExists{xurl.sty}{\usepackage{xurl}}{} % add URL line breaks if available
\IfFileExists{bookmark.sty}{\usepackage{bookmark}}{\usepackage{hyperref}}
\hypersetup{
  pdftitle={Manual de Gestão de Dados de Recursos do Solo},
  pdfauthor={Alessandro Samuel-Rosa},
  colorlinks=true,
  linkcolor=Maroon,
  filecolor=Maroon,
  citecolor=Blue,
  urlcolor=Blue,
  pdfcreator={LaTeX via pandoc}}
\urlstyle{same} % disable monospaced font for URLs
\usepackage{longtable,booktabs}
% Correct order of tables after \paragraph or \subparagraph
\usepackage{etoolbox}
\makeatletter
\patchcmd\longtable{\par}{\if@noskipsec\mbox{}\fi\par}{}{}
\makeatother
% Allow footnotes in longtable head/foot
\IfFileExists{footnotehyper.sty}{\usepackage{footnotehyper}}{\usepackage{footnote}}
\makesavenoteenv{longtable}
\usepackage{graphicx,grffile}
\makeatletter
\def\maxwidth{\ifdim\Gin@nat@width>\linewidth\linewidth\else\Gin@nat@width\fi}
\def\maxheight{\ifdim\Gin@nat@height>\textheight\textheight\else\Gin@nat@height\fi}
\makeatother
% Scale images if necessary, so that they will not overflow the page
% margins by default, and it is still possible to overwrite the defaults
% using explicit options in \includegraphics[width, height, ...]{}
\setkeys{Gin}{width=\maxwidth,height=\maxheight,keepaspectratio}
% Set default figure placement to htbp
\makeatletter
\def\fps@figure{htbp}
\makeatother
\setlength{\emergencystretch}{3em} % prevent overfull lines
\providecommand{\tightlist}{%
  \setlength{\itemsep}{0pt}\setlength{\parskip}{0pt}}
\setcounter{secnumdepth}{5}
\usepackage{booktabs}
\usepackage[utf8]{inputenc}
\usepackage[T1]{fontenc}
\usepackage[brazilian]{babel}
% https://www.ctan.org/pkg/fancyhdr
\usepackage{fancyhdr}
\pagestyle{fancy}
\fancyhead[RO]{\textsl{\nouppercase{\rightmark}}}
\fancyhead[LE]{\textsl{\nouppercase{\leftmark}}}
\fancyfoot[LE, RO]{\thepage}
\renewcommand{\headrulewidth}{0.4pt}
% http://tug.ctan.org/tex-archive/macros/latex/contrib/tufte-latex/
\publisher{Universidade Tecnológica Federal do Paraná}
% https://tex.stackexchange.com/questions/443018/package-inputenc-error-unicode-char-%CC%81-u301inputenc
% \DeclareUnicodeCharacter{0301}{*************************************}
% \DeclareUnicodeCharacter{2265}{AAAAAAAAAAAAAAAAAAAAAAAAAAAAAAAAAAAAA}
% https://tex.stackexchange.com/questions/473859/subtitle-in-tufte-book
\makeatletter
\newcommand{\plainsubtitle}{}%     plain-text-only subtitle
\newcommand{\subtitle}[1]{%
  \gdef\@subtitle{#1}%
  \renewcommand{\plainsubtitle}{#1}% use provided plain-text title
  \ifthenelse{\isundefined{\hypersetup}}%
    {}% hyperref is not loaded; do nothing
    {\hypersetup{pdftitle={\plaintitle: \plainsubtitle{}}}}% set the PDF metadata title
}
\renewcommand{\maketitlepage}[0]{%
  \cleardoublepage%
  {%
  \sffamily%
  \begin{fullwidth}%
    \fontsize{18}{20}\selectfont\par\noindent\textcolor{darkgray}{\allcaps{\thanklessauthor}}%
    \vspace{11.5pc}%
    \fontsize{36}{40}\selectfont\par\noindent\textcolor{darkgray}{\allcaps{\thanklesstitle}}%
    \vspace{5pc}%
    \fontsize{24}{28}\selectfont\par\noindent\textcolor{darkgray}{\allcaps{\plainsubtitle}}%
    \vfill%
    \fontsize{14}{16}\selectfont\par\noindent\allcaps{\thanklesspublisher}%
  \end{fullwidth}%
  }
  \thispagestyle{empty}%
  \clearpage%
}
\makeatother
% https://tex.stackexchange.com/questions/77999/remove-indent-of-paragraph-and-add-line-skip-with-tufte-latex/
\makeatletter
% Paragraph indentation and separation for normal text
\renewcommand{\@tufte@reset@par}{%
  \setlength{\RaggedRightParindent}{0.0pc}%
  \setlength{\JustifyingParindent}{0.0pc}%
  \setlength{\parindent}{0pc}%
  % \setlength{\parskip}{\baselineskip}%
}
\@tufte@reset@par
\pagecolor{Emerald}

\title{Manual de Gestão de Dados de Recursos do Solo}
\author{Alessandro Samuel-Rosa}
\date{2020-12-29}

\begin{document}
\maketitle

\newpage
\nopagecolor
\begin{fullwidth}
~\vfill
\thispagestyle{empty}
\setlength{\parindent}{0pt}
\setlength{\parskip}{\baselineskip}

\par{\smallcaps{Publicado pela Universidade Tecnológica Federal do Paraná, Câmpus Santa Helena}}

\par{
  Curso de Agronomia/Laboratório de Pedometria\\
  Prolongamento da Rua Cerejeira, s/n\\
  Bairro São Luiz\\
  CEP 85892-000\\
  Santa Helena - PR - Brasil\\
  Telefone Geral +55 (45) 3268-8800
}

\par\smallcaps{www.pedometria.org}

\par\includegraphics[width=1.5cm]{img/cc-by-nc-sa.png} Exceto quando proveniente de outras fontes ou onde especificado o contrário, o conteúdo deste documento está licenciado sob uma licença \href{http://creativecommons.org/licenses/by-nc-sa/4.0/}{Creative Commons Atribuição-NãoComercial-CompartilhaIgual 4.0 Internacional}.

\par\textit{Última edição em \today.}
\end{fullwidth}

{
\hypersetup{linkcolor=}
\setcounter{tocdepth}{1}
\tableofcontents
}
\hypertarget{prefuxe1cio}{%
\chapter*{Prefácio}\label{prefuxe1cio}}
\addcontentsline{toc}{chapter}{Prefácio}

\hypertarget{intro}{%
\chapter{Introdução}\label{intro}}

\hypertarget{variuxe1veis-morfoluxf3gicas}{%
\chapter*{Variáveis Morfológicas}\label{variuxe1veis-morfoluxf3gicas}}
\addcontentsline{toc}{chapter}{Variáveis Morfológicas}

\hypertarget{cor-do-solo}{%
\chapter{Cor do Solo}\label{cor-do-solo}}

\emph{Alessandro Samuel-Rosa}\footnote{Dr.~Alessandro Samuel-Rosa é professor na Universidade Tecnológica Federal do Paraná, Departamento de Agronomia, Prolongamento da Rua Cerejeira, s/n, CEP 85892-000, Santa Helena, Paraná, Brasil.}, \emph{Taciara Zborowski Horst-Heinen}\footnote{MSc. Taciara Zborowski Horst-Heinen é pesquisadora na Universidade Federal de Santa Maria, Programa de Pós-Graduação em Ciência do Solo, Avenida Roraima, 1000, CEP 97105-900, Santa Maria, Rio Grande do Sul, Brasil.}

\hypertarget{conceitos-e-definiuxe7uxf5es}{%
\section{Conceitos e definições}\label{conceitos-e-definiuxe7uxf5es}}

A \textbf{cor do solo} é uma qualidade relativa à percepção visual que seres humanos têm do solo. Essa percepção é resultado da impressão que a luz refletida pelo solo produz nos olhos de quem o está observando. Logo, trata-se de uma qualidade perceptível apenas na presença de luz visível, seja ela proveniente de fontes artificiais ou naturais.

A luz visível corresponde aos comprimentos de onda entre cerca de 400 e 700 nm. Essa porção do espectro eletromagnético é popularmente conhecida como espectro óptico ou espectro do visível. Se ignorarmos os principais fatores intervenientes como imperfeições do olho ou cérebro do observador, condições do ambiente circunvizinho, intensidade da luz, e ângulo de incidência e visada, a cor do solo resulta da repartição diferencial dos comprimentos de onda da luz visível incidente entre os fenômenos de absorção, transmissão e reflexão. Essa repartição diferencial é fruto da maneira como os vários constituíntes do solo interagem com a radiação eletromagnética. É a radiação presente nos comprimentos de onda do visível (fótons) que deixam a superfície do solo em maior quantidade, por reflexão ou, inclusive, emissão, que ditam a cor do solo.

\hypertarget{muxe9todos-de-determinauxe7uxe3o}{%
\section{Métodos de determinação}\label{muxe9todos-de-determinauxe7uxe3o}}

\hypertarget{fontes-de-variauxe7uxe3o}{%
\section{Fontes de variação}\label{fontes-de-variauxe7uxe3o}}

\hypertarget{padronizauxe7uxe3o}{%
\section{Padronização}\label{padronizauxe7uxe3o}}

O esquema padronizado de identificação e descrição de métodos de determinação da cor de materiais do solo utilizado no FEBR é apresentado na tabela abaixo. São considerados dois materiais de solo (matriz e manchas), dois estados de umidade (úmido e seco) e duas condições mecânicas (amassada e triturada).

\begin{longtable}[]{@{}ll@{}}
\toprule
\begin{minipage}[b]{0.51\columnwidth}\raggedright
Código de identificação\strut
\end{minipage} & \begin{minipage}[b]{0.43\columnwidth}\raggedright
Descrição mínima sugerida\strut
\end{minipage}\tabularnewline
\midrule
\endhead
\begin{minipage}[t]{0.51\columnwidth}\raggedright
\texttt{cor\_matriz\_umido\_munsell}\strut
\end{minipage} & \begin{minipage}[t]{0.43\columnwidth}\raggedright
Notação da cor de material de solo {[}\texttt{cor}{]}. Determinada na matriz do solo {[}\texttt{matriz}{]}. Amostra em estado de umidade úmido {[}\texttt{umido}{]}. Expressa utilizando a notação de Munsell® (matiz valor/croma) {[}\texttt{munsell}{]}. \{A DETERMINAÇÃO FOI FEITA NO CAMPO OU EM LABORATÓRIO?\}\strut
\end{minipage}\tabularnewline
\begin{minipage}[t]{0.51\columnwidth}\raggedright
\texttt{cor\_matriz\_seco\_munsell}\strut
\end{minipage} & \begin{minipage}[t]{0.43\columnwidth}\raggedright
Notação da cor de material de solo {[}\texttt{cor}{]}. Determinada na matriz do solo {[}\texttt{matriz}{]}. Amostra em estado de umidade seco {[}seco{]}. Expressa utilizando a notação de Munsell® (matiz valor/croma) {[}\texttt{munsell}{]}. \{ESPECIFICAR SE A DETERMINAÇÃO FOI FEITA NO CAMPO OU EM LABORATÓRIO\}\strut
\end{minipage}\tabularnewline
\begin{minipage}[t]{0.51\columnwidth}\raggedright
\texttt{cor\_matriz\_amassada\_munsell}\strut
\end{minipage} & \begin{minipage}[t]{0.43\columnwidth}\raggedright
Notação da cor de material de solo {[}\texttt{cor}{]}. Determinada na matriz do solo {[}\texttt{matriz}{]}. Amostra em estado de umidade úmido e condição mecânica amassada {[}\texttt{amassada}{]}. Expressa utilizando a notação de Munsell® (matiz valor/croma) {[}\texttt{munsell}{]}. \{ESPECIFICAR SE A DETERMINAÇÃO FOI FEITA NO CAMPO OU EM LABORATÓRIO\}\strut
\end{minipage}\tabularnewline
\begin{minipage}[t]{0.51\columnwidth}\raggedright
\texttt{cor\_matriz\_triturada\_munsell}\strut
\end{minipage} & \begin{minipage}[t]{0.43\columnwidth}\raggedright
Notação da cor de material de solo {[}\texttt{cor}{]}. Determinada na matriz do solo {[}\texttt{matriz}{]}. Amostra em estado de umidade seco e condição mecânica triturada {[}\texttt{triturada}{]}. Expressa utilizando a notação de Munsell® (matiz valor/croma) {[}\texttt{munsell}{]}. \{ESPECIFICAR SE A DETERMINAÇÃO FOI FEITA NO CAMPO OU EM LABORATÓRIO\}\strut
\end{minipage}\tabularnewline
\begin{minipage}[t]{0.51\columnwidth}\raggedright
\texttt{cor\_mancha\_umido\_munsell}\strut
\end{minipage} & \begin{minipage}[t]{0.43\columnwidth}\raggedright
Notação da cor de material de solo {[}\texttt{cor}{]}. Determinada na mancha do solo {[}\texttt{mancha}{]}. Amostra em estado de umidade úmido {[}\texttt{umido}{]}. Expressa utilizando a notação de Munsell® (matiz valor/croma) {[}\texttt{munsell}{]}. \{ESPECIFICAR SE A DETERMINAÇÃO FOI FEITA NO CAMPO OU EM LABORATÓRIO\}\strut
\end{minipage}\tabularnewline
\bottomrule
\end{longtable}

\hypertarget{exemplo}{%
\section{Exemplo}\label{exemplo}}

A tabela a seguir apresenta um exemplo de organização tabular de dados de cor da matriz e manchas de um perfil de solo. Os dados são reais e foram obtidos da descrição morfológica de um perfil de solo classificado como Planossolo Nátrico Órtico vertissólico, localizado no município de Pacaraima, Roraima\footnote{Virlei Álvaro de Oliveira et al., ``Capítulo 2. Solos Da Xi Reunião Brasileira de Classificação E Correlação de Solos (Rcc de Roraima),'' in \emph{Guia de Campo Da Xi Reunião Brasileira de Classificação E Correlação de Solos: RCC de Roraima}, First (Brasília, DF: EMBRAPA, 2018), 12--29.}. A tabela está transposta, com as observações nas colunas e as variáveis nas linhas, para facilitar a visualização e comparação dos códigos de identificação da cores da matriz e machas do perfil de solo.

Dentre os oito horizontes, quatro possuem matriz com duas ou mais cores (Btgn, Btgnv1, Btgnv2 e 2Cgn). Nestes, as cores da matriz ocupam áreas superficiais aproximadamente similares, com distribuição espacial muito complexa, que dificulta a identificação visual de uma cor dominante com credibilidade. Trata-se do padrão multicolorido intricado conhecido como variegado.

\begin{longtable}[]{@{}lllllllll@{}}
\toprule
\begin{minipage}[b]{0.22\columnwidth}\raggedright
\strut
\end{minipage} & \begin{minipage}[b]{0.06\columnwidth}\raggedright
An\strut
\end{minipage} & \begin{minipage}[b]{0.06\columnwidth}\raggedright
En\strut
\end{minipage} & \begin{minipage}[b]{0.07\columnwidth}\raggedright
EBn\strut
\end{minipage} & \begin{minipage}[b]{0.07\columnwidth}\raggedright
Btn\strut
\end{minipage} & \begin{minipage}[b]{0.07\columnwidth}\raggedright
Btgn\strut
\end{minipage} & \begin{minipage}[b]{0.07\columnwidth}\raggedright
Btgnv1\strut
\end{minipage} & \begin{minipage}[b]{0.07\columnwidth}\raggedright
Btgnv2\strut
\end{minipage} & \begin{minipage}[b]{0.06\columnwidth}\raggedright
2Cgn\strut
\end{minipage}\tabularnewline
\midrule
\endhead
\begin{minipage}[t]{0.22\columnwidth}\raggedright
\texttt{cor\_matriz\_umido\_munsell\_1}\strut
\end{minipage} & \begin{minipage}[t]{0.06\columnwidth}\raggedright
10YR 5/1\strut
\end{minipage} & \begin{minipage}[t]{0.06\columnwidth}\raggedright
10YR 5/1\strut
\end{minipage} & \begin{minipage}[t]{0.07\columnwidth}\raggedright
10YR 6/2\strut
\end{minipage} & \begin{minipage}[t]{0.07\columnwidth}\raggedright
7,5YR 5/2\strut
\end{minipage} & \begin{minipage}[t]{0.07\columnwidth}\raggedright
7,5YR 6/2\strut
\end{minipage} & \begin{minipage}[t]{0.07\columnwidth}\raggedright
10YR 4/1\strut
\end{minipage} & \begin{minipage}[t]{0.07\columnwidth}\raggedright
10YR 4/1\strut
\end{minipage} & \begin{minipage}[t]{0.06\columnwidth}\raggedright
10YR 6/2\strut
\end{minipage}\tabularnewline
\begin{minipage}[t]{0.22\columnwidth}\raggedright
\texttt{cor\_matriz\_umido\_munsell\_2}\strut
\end{minipage} & \begin{minipage}[t]{0.06\columnwidth}\raggedright
-\strut
\end{minipage} & \begin{minipage}[t]{0.06\columnwidth}\raggedright
-\strut
\end{minipage} & \begin{minipage}[t]{0.07\columnwidth}\raggedright
-\strut
\end{minipage} & \begin{minipage}[t]{0.07\columnwidth}\raggedright
-\strut
\end{minipage} & \begin{minipage}[t]{0.07\columnwidth}\raggedright
7,5YR 5/2\strut
\end{minipage} & \begin{minipage}[t]{0.07\columnwidth}\raggedright
10YR 5/2\strut
\end{minipage} & \begin{minipage}[t]{0.07\columnwidth}\raggedright
10YR 5/2\strut
\end{minipage} & \begin{minipage}[t]{0.06\columnwidth}\raggedright
10YR 7/2\strut
\end{minipage}\tabularnewline
\begin{minipage}[t]{0.22\columnwidth}\raggedright
\texttt{cor\_matriz\_umido\_munsell\_3}\strut
\end{minipage} & \begin{minipage}[t]{0.06\columnwidth}\raggedright
-\strut
\end{minipage} & \begin{minipage}[t]{0.06\columnwidth}\raggedright
-\strut
\end{minipage} & \begin{minipage}[t]{0.07\columnwidth}\raggedright
-\strut
\end{minipage} & \begin{minipage}[t]{0.07\columnwidth}\raggedright
-\strut
\end{minipage} & \begin{minipage}[t]{0.07\columnwidth}\raggedright
-\strut
\end{minipage} & \begin{minipage}[t]{0.07\columnwidth}\raggedright
10YR 6/2\strut
\end{minipage} & \begin{minipage}[t]{0.07\columnwidth}\raggedright
10YR 4/2\strut
\end{minipage} & \begin{minipage}[t]{0.06\columnwidth}\raggedright
-\strut
\end{minipage}\tabularnewline
\begin{minipage}[t]{0.22\columnwidth}\raggedright
\texttt{cor\_matriz\_umido\_munsell\_4}\strut
\end{minipage} & \begin{minipage}[t]{0.06\columnwidth}\raggedright
-\strut
\end{minipage} & \begin{minipage}[t]{0.06\columnwidth}\raggedright
-\strut
\end{minipage} & \begin{minipage}[t]{0.07\columnwidth}\raggedright
-\strut
\end{minipage} & \begin{minipage}[t]{0.07\columnwidth}\raggedright
-\strut
\end{minipage} & \begin{minipage}[t]{0.07\columnwidth}\raggedright
-\strut
\end{minipage} & \begin{minipage}[t]{0.07\columnwidth}\raggedright
7,5YR 5/8\strut
\end{minipage} & \begin{minipage}[t]{0.07\columnwidth}\raggedright
7,5YR 5/6\strut
\end{minipage} & \begin{minipage}[t]{0.06\columnwidth}\raggedright
-\strut
\end{minipage}\tabularnewline
\begin{minipage}[t]{0.22\columnwidth}\raggedright
\texttt{cor\_matriz\_seco\_munsell}\strut
\end{minipage} & \begin{minipage}[t]{0.06\columnwidth}\raggedright
10YR 6/2\strut
\end{minipage} & \begin{minipage}[t]{0.06\columnwidth}\raggedright
10YR 6/1\strut
\end{minipage} & \begin{minipage}[t]{0.07\columnwidth}\raggedright
-\strut
\end{minipage} & \begin{minipage}[t]{0.07\columnwidth}\raggedright
-\strut
\end{minipage} & \begin{minipage}[t]{0.07\columnwidth}\raggedright
-\strut
\end{minipage} & \begin{minipage}[t]{0.07\columnwidth}\raggedright
-\strut
\end{minipage} & \begin{minipage}[t]{0.07\columnwidth}\raggedright
-\strut
\end{minipage} & \begin{minipage}[t]{0.06\columnwidth}\raggedright
-\strut
\end{minipage}\tabularnewline
\begin{minipage}[t]{0.22\columnwidth}\raggedright
\texttt{cor\_mancha\_umido\_munsell\_1}\strut
\end{minipage} & \begin{minipage}[t]{0.06\columnwidth}\raggedright
-\strut
\end{minipage} & \begin{minipage}[t]{0.06\columnwidth}\raggedright
-\strut
\end{minipage} & \begin{minipage}[t]{0.07\columnwidth}\raggedright
7,5YR 5/8\strut
\end{minipage} & \begin{minipage}[t]{0.07\columnwidth}\raggedright
7,5YR 5/8\strut
\end{minipage} & \begin{minipage}[t]{0.07\columnwidth}\raggedright
-\strut
\end{minipage} & \begin{minipage}[t]{0.07\columnwidth}\raggedright
-\strut
\end{minipage} & \begin{minipage}[t]{0.07\columnwidth}\raggedright
-\strut
\end{minipage} & \begin{minipage}[t]{0.06\columnwidth}\raggedright
10YR 6/8\strut
\end{minipage}\tabularnewline
\begin{minipage}[t]{0.22\columnwidth}\raggedright
\texttt{cor\_mancha\_umido\_munsell\_2}\strut
\end{minipage} & \begin{minipage}[t]{0.06\columnwidth}\raggedright
-\strut
\end{minipage} & \begin{minipage}[t]{0.06\columnwidth}\raggedright
-\strut
\end{minipage} & \begin{minipage}[t]{0.07\columnwidth}\raggedright
-\strut
\end{minipage} & \begin{minipage}[t]{0.07\columnwidth}\raggedright
7,5YR 4/1\strut
\end{minipage} & \begin{minipage}[t]{0.07\columnwidth}\raggedright
-\strut
\end{minipage} & \begin{minipage}[t]{0.07\columnwidth}\raggedright
-\strut
\end{minipage} & \begin{minipage}[t]{0.07\columnwidth}\raggedright
-\strut
\end{minipage} & \begin{minipage}[t]{0.06\columnwidth}\raggedright
-\strut
\end{minipage}\tabularnewline
\begin{minipage}[t]{0.22\columnwidth}\raggedright
\texttt{cor\_mancha\_umido\_munsell\_3}\strut
\end{minipage} & \begin{minipage}[t]{0.06\columnwidth}\raggedright
-\strut
\end{minipage} & \begin{minipage}[t]{0.06\columnwidth}\raggedright
-\strut
\end{minipage} & \begin{minipage}[t]{0.07\columnwidth}\raggedright
-\strut
\end{minipage} & \begin{minipage}[t]{0.07\columnwidth}\raggedright
7,5YR 7/2\strut
\end{minipage} & \begin{minipage}[t]{0.07\columnwidth}\raggedright
-\strut
\end{minipage} & \begin{minipage}[t]{0.07\columnwidth}\raggedright
-\strut
\end{minipage} & \begin{minipage}[t]{0.07\columnwidth}\raggedright
-\strut
\end{minipage} & \begin{minipage}[t]{0.06\columnwidth}\raggedright
-\strut
\end{minipage}\tabularnewline
\bottomrule
\end{longtable}

\hypertarget{variuxe1veis-fuxedsicas}{%
\chapter*{Variáveis Físicas}\label{variuxe1veis-fuxedsicas}}
\addcontentsline{toc}{chapter}{Variáveis Físicas}

\hypertarget{densidade}{%
\chapter{Densidade}\label{densidade}}

\emph{Alessandro Samuel-Rosa}\footnote{Dr.~Alessandro Samuel-Rosa é professor na Universidade Tecnológica Federal do Paraná, Departamento de Agronomia, Prolongamento da Rua Cerejeira, s/n, CEP 85892-000, Santa Helena, Paraná, Brasil.}, \emph{Wenceslau Geraldes Teixeira}\footnote{Dr.~Wenceslau Geraldes Teixeira é pesquisador na Empresa Brasileira de Pesquisa Agropecuária, Centro Nacional de Pesquisa de Solos, Rua Jardim Botânico, 1024, CEP 22460-000, Rio de Janeiro, Rio de Janeiro, Brasil.}, \emph{João Herbert Moreira Viana}\footnote{Dr.~João Herbert Moreira Viana é pesquisador na Empresa Brasileira de Pesquisa Agropecuária, Centro Nacional de Pesquisa de Milho e Sorgo, Rodovia MG-424, Km 45, CEP 35701-970, Sete Lagoas, Minas Gerais, Brasil.}

\hypertarget{conceitos-e-definiuxe7uxf5es-1}{%
\section{Conceitos e definições}\label{conceitos-e-definiuxe7uxf5es-1}}

A densidade do solo, também chamada erroneamente de densidade aparente ou densidade global, é uma medida da massa seca de todo o material contido em determinado volume de solo. Esse material inclui todas as frações dos componentes de origem mineral e orgânica do solo. No FEBR, a densidade do solo é definida, operacionalmente, como densidade do solo inteiro, e seu código de identificação é \texttt{densidade}. No World Soil Information Service (WoSIS), a densidade do solo inteiro é denominada \emph{bulk density whole soil}, e seu código de identificação é BDWS.

\hypertarget{muxe9todos-de-determinauxe7uxe3o-1}{%
\section{Métodos de determinação}\label{muxe9todos-de-determinauxe7uxe3o-1}}

A DSI pode ser determinada utilizando uma variedade de métodos diretos e indiretos. Aqui tratamos apenas dos métodos diretos. Esses métodos diferenciam-se em si pela maneira de obtenção das amostras do solo e quantificação (ou estimativa) da sua massa e volume. Em geral, a escolha por um ou outro método está relacionada às características do solo, especialmente a granulometria e composição mineralógica do solo. Os principais métodos diretos de determinação da DSI são: cilindro, torrão, proveta, monolito, e escavação. Assim, o segundo nível de codificação da DSI no FEBR consiste na identificação desses métodos, ou seja:

\begin{itemize}
\tightlist
\item
  \texttt{densidade\_cilindro},
\item
  \texttt{densidade\_torrao},
\item
  \texttt{densidade\_proveta},
\item
  \texttt{densidade\_monolito}, e
\item
  \texttt{densidade\_escavacao}.
\end{itemize}

Dentre os métodos diretos, \texttt{densidade\_cilindro} e \texttt{densidade\_torrao} são os mais utilizados quando o material do solo é coletado tentando se preservar a estrutura do solo. Também costumam ser os mais utilizados quando o solo é composto, em sua quase totalidade, pela fração fina (\textgreater{} 2 mm) e, no caso de \texttt{densidade\_torrao}, quando há predomínio da fração argila. Já o método \texttt{densidade\_proveta} aplica-se mais aos casos em que há dominância da fração areia ou quando se pretende fazer estimativas da DSI usando apenas a fração fina do solo. Por fim, os métodos \texttt{densidade\_monolito} e \texttt{densidade\_escavacao} são mais usados quando há predomínio das frações grossas (\textgreater{} 2 mm), que impossibilitam ou dificultam sobremaneira a inserção de cilindros no solo ou coleta de torrões com estrutura preservada. As mesmas dificuldades são encontradas em solo com expressivo volume de raízes de grande diâmetro e elevado conteúdo de matéria orgânica, como em florestas, onde o método \texttt{densidade\_escavacao} é o mais usado.

\hypertarget{fontes-de-variauxe7uxe3o-1}{%
\section{Fontes de variação}\label{fontes-de-variauxe7uxe3o-1}}

Cada um dos métodos diretos de determinação da DSI apontados acima abrange uma variedade de detalhes metodológicos que podem exercer pequenas influências sobre os resultados obtidos. Por exemplo, no caso do \texttt{densidade\_cilindro}, há chance de influência do método utilizado para inserção do cilindro no solo, que pode ser por percussão com martelo ou por pressão com prensa hidráulica. Em geral, espera-se que o método da percussão com martelo cause maior compactação e estilhaçamento da amostra, especialmente quando o cilindro usado é muito pequeno.\footnote{Manuel Casanova et al., ``Direct Measurement and Prediction of Bulk Density on Alluvial Soils of Central Chile,'' \emph{Chilean Journal of Agricultural Research} 76, no. 1 (March 2016): 105--13, \url{https://doi.org/10.4067/s0718-58392016000100015}.} A compactação da amostra também pode ocorrer quando os cilindros usados não possuem corte em bisel na face inserida no solo. Por outro lado, espera-se que, quanto mais finas as paredes do cilindro, menor seja a compactação da amostra.

Outro detalhe metodológico que pode influenciar os resultados é o tamanho do cilindro.\footnote{Ahmed Abed Gatea Al-Shammary et al., ``Soil Bulk Density Estimation Methods: A Review,'' \emph{Pedosphere} 28, no. 4 (August 2018): 581--96, \url{https://doi.org/10.1016/s1002-0160(18)60034-7}; Casanova et al., ``Direct Measurement and Prediction of Bulk Density on Alluvial Soils of Central Chile.''} Boa parte dos estudos usa cilindros de 100 cm3, tamanho recomendado pelo Manual de Métodos de Análise de Solo.\footnote{Brivaldo Gomes de Almeida et al., ``Densidade Do Solo,'' in \emph{Manual de Métodos de Análise de Solo}, ed. Paulo César Teixeira et al. (Brasília: EMBRAPA, 2017), 65--75.} Contudo, relatos do uso de cilindros de diversos tamanhos são encontrados na literatura, por exemplo, 9,65 cm3 (amostras pequenas usadas em reometria), 30 cm3, 300 cm3, e 1000 cm3 (calibração de sensores de medição do conteúdo de água no solo).

A definição do tamanho do cilindro é importante pois está relacionada ao volume elementar representativo (REV, do inglês \emph{representative elementar volume}) de cada tipo de solo. Por exemplo, se há poros muito grandes ou cascalhos e calhaus no sistema, cilindros pequenos não seriam representativos. A questão do REV se aplica também aos demais métodos, principalmente \texttt{densidade\_torrao}, \texttt{densidade\_monolito} e \texttt{densidade\_escavacao}. Especialmente no caso do \texttt{densidade\_monolito} e \texttt{densidade\_escavacao}, que utilizam amostras maiores, a questão da REV seria melhor atendida. Esse já não seria o caso do método \texttt{densidade\_torrao}, onde torrões de tamanhos bastantes variados são usados, mas geralmente menores que 100 m³ devido à dificuldade de se obter agregados estáveis -- especialmente em solo de granulometria mais grossa. Um aspecto importante a ser considerado aqui é que o uso de repetições locais não resolve o problema da REV. Isso porque a função das repetições é quantificar o erro de medida, não resolver variação espacial de curtíssima distância -- no FEBR, repetições são identificada na tabela \emph{camada} usando o campo \texttt{amostra\_id}.

Diferentes dos demais métodos, \texttt{densidade\_torrao} e \texttt{densidade\_monolito} incluem a impermeabilização das amostras. Esse procedimento pode ser feito usando diferentes compostos, tais como parafina, verniz ou querosene. Apesar das diferenças entre esses compostos, espera-se que o uso de diferentes impermeabilizantes tenha efeito muitíssimo pequeno, se algum, sobre os resultados de DSI obtidos. O mesmo pode ser esperado para o caso do \texttt{densidade\_escavacao}, método em que podem ser utilizados diferentes materiais para preenchimento da escavação, por exemplo, areia ou água.

Um detalhe metodológico de grande importância na determinação da DSI é a umidade da amostra do solo no momento da quantificação da sua massa e volume. Na maioria dos casos, esses valores são obtidos com a amostra seca. Contudo, para as amostras de alguns tipos de solo, o que se faz é equilibrar as amostras em determinado potencial, geralmente 1/3 atm ou 33 kPa, que corresponde ao que se convencionou chamar de capacidade de campo. Esse procedimento se justifica nos casos de solos com presença de argilas expansivas em quantidade suficiente para causar aumento perceptível do volume da massa de solo, como no caso dos solos com características vérticas e Vertissolos. Na determinação da DSI desses solos, o volume e massa da amostra são quantificados nesta umidade controlada.\footnote{Clément Mathieu and Françoise Pieltain, \emph{Analyse Physique Des Sols: Méthodes Choisies} (Paris: Lavoisier, 1998); R. B. Grossman and T. G. Reinsch, ``2.1 Bulk Density and Linear Extensibility,'' in \emph{SSSA Book Series}, ed. Jacob H. Dane and G. Clarke Topp (Madison, WI, USA: Soil Science Society of America, 2002), 201--28, \url{https://doi.org/10.2136/sssabookser5.4.c9}.}

Dentre todos os detalhes metodológicos mencionados acima, a condição de umidade da amostra de solo no momento da quantificação da massa e volume é aquele com maior potencial de exercer influências sobre os resultados obtidos. Para a maioria das outras fontes adicionais de variação nos resultados, os trabalhos de comparação, quando existentes, são de difícil acesso. No FEBR, ainda não foram registrados conjuntos de dados em que a determinação da DSI tenha sido realizada com amostras úmidas, razão pela qual não se faz necessário definir codificação adicional. Contudo, caso dados dessa natureza sejam submetidos ao FEBR, a codificação poderá ser definida, em seu terceiro nível, pela condição da amostra no momento da quantificação da massa e volume -- por exemplo, \texttt{estufa}, \texttt{ar}, \texttt{campo}, \texttt{33kpa}, entre outras. Esse procedimento estaria em acordo com aqueles adotados no World Soil Information System (WoSIS).\footnote{Eloi Ribeiro, Niels H. Batjes, and A. J. M. van Oostrum, \emph{World Soil Information Service (Wosis) - Towards the Standardization and Harmonization of World Soil Data. Procedures Manual 2018} (Wageningen, The Netherlands: ISRIC World Soil Information, 2018), \url{https://doi.org/10.17027/isric-wdcsoils.20180001}.}

\hypertarget{padronizauxe7uxe3o-1}{%
\section{Padronização}\label{padronizauxe7uxe3o-1}}

A Tabela 1 apresenta os códigos padronizados utilizados no FEBR para identificar cada um dos métodos de determinação da densidade do solo inteiro e da fração terra fina. Também apresenta uma descrição mínima sugerida para cada um desses métodos, ou seja, um roteiro padronizado para especificação dos detalhes mais importantes de cada método. Detalhes adicionais dos métodos podem ser especificados conforme indicado pelo texto em caixa alta entre chaves.

Tabela 1. Codificação e descrição mínima sugerida da densidade do solo inteiro e da fração terra fina no FEBR.

A Tabela 2 abaixo apresenta as especificações de unidade de medida, a precisão numérica, tipo de dado, e categoria da variável no FEBR.

Tabela 2. Especificações da densidade do solo inteiro no FEBR.

\hypertarget{harmonizauxe7uxe3o}{%
\section{Harmonização}\label{harmonizauxe7uxe3o}}

Existe uma grande variedade de métodos de determinação direta da DSI. Cada um deles foi desenvolvido de maneira a permitir a determinação da DSI em diferentes condições, por exemplo, de granulometria e mineralogia do solo. Assim, não existe um único método de referência para a determinação da DSI. Isso significa que, em geral, não é possível transformar os valores obtidos com um método qualquer para valores aproximadamente equivalentes aos valores que seriam obtidos com outro método. Dados empíricos que permitam construir modelos estatísticos para fazer essa conversão são escassos e pouco representativos.

A limitação de dados empíricos impossibilita, por ora, a harmonização dos dados de DSI no FEBR. Estudos que venham resolver essa limitação são bem vindos, especialmente no caso dos métodos que guardam mais similaridades entre si, por exemplo, \texttt{densidade\_torrao} e \texttt{densidade\_cilindro}. Em geral, \texttt{densidade\_torrao} tende a sobre estimar a DSI, comparado ao \texttt{densidade\_cilindro}, por desconsiderar o espaço poroso existente entre os agregados.\footnote{Casanova et al., ``Direct Measurement and Prediction of Bulk Density on Alluvial Soils of Central Chile.''} O FEBR tem condições de acomodar um modelo estatístico que transforme os valores de DSI obtidos com \texttt{densidade\_torrao} para valores aproximadamente equivalentes àqueles que seriam obtidos com \texttt{densidade\_cilindro} -- e vice-versa.

\hypertarget{agradecimentos}{%
\section{Agradecimentos}\label{agradecimentos}}

\begin{itemize}
\tightlist
\item
  Iêde de Brito Chaves, professor aposentado da Universidade Federal da Paraíba (UFPB), pelas considerações acerca das potenciais dificuldades causadas pelo uso das expressões `densidade do solo inteiro', `densidade do solo', `densidade aparente' e `densidade global'. {[}\href{https://groups.google.com/d/msgid/soil-mapping/7c8745ca-243d-4d15-97cc-2d00cb5eefdd\%40googlegroups.com}{Ver essa discussão na Web}{]}
\item
  Anderson Sandro da Rocha, professor da Universidade Tecnológica Federal do Paraná, Curso de Agronomia.
\item
  Paulo Ivonir Gubiani, professor da Universidade Federal de Santa Maria, Departamento de Solos.
\end{itemize}

\hypertarget{variuxe1veis-espectrais}{%
\chapter*{Variáveis Espectrais}\label{variuxe1veis-espectrais}}
\addcontentsline{toc}{chapter}{Variáveis Espectrais}

\hypertarget{reflectuxe2ncia}{%
\chapter{Reflectância}\label{reflectuxe2ncia}}

\emph{Taciara Zborowski Horst-Heinen}\footnote{MSc. Taciara Zborowski Horst-Heinen é pesquisadora na Universidade Federal de Santa Maria, Programa de Pós-Graduação em Ciência do Solo, Avenida Roraima, 1000, CEP 97105-900, Santa Maria, Rio Grande do Sul, Brasil.}, \emph{Alessandro Samuel-Rosa}\footnote{Dr.~Alessandro Samuel-Rosa é professor na Universidade Tecnológica Federal do Paraná, Departamento de Agronomia, Prolongamento da Rua Cerejeira, s/n, CEP 85892-000, Santa Helena, Paraná, Brasil.}

\hypertarget{conceitos-e-definiuxe7uxf5es-2}{%
\section{Conceitos e definições}\label{conceitos-e-definiuxe7uxf5es-2}}

\hypertarget{muxe9todos-de-determinauxe7uxe3o-2}{%
\section{Métodos de determinação}\label{muxe9todos-de-determinauxe7uxe3o-2}}

\footnote{R. A. Viscarra Rossel et al., ``A Global Spectral Library to Characterize the World's Soil,'' \emph{Earth-Science Reviews} 155 (2016): 198--230, \url{https://doi.org/10.1016/j.earscirev.2016.01.012}.}

\hypertarget{padronizauxe7uxe3o-2}{%
\section{Padronização}\label{padronizauxe7uxe3o-2}}

W350\ldots W2500Wavebands in nanometers at which reflectance values are given in the database.\footnote{ICRAF and ISRIC, ``A Globally Distributed Soil Spectral Library Visible Near Infrared Diffuse Reflectance Spectra'' (Nairobi, Kenya: World Agroforestry Centre (ICRAF), 2010).}

\begin{tabular}{l|l}
\hline
campo\_id & campo\_descricao\\
\hline
`reflectancia\_seco\_lab` & Reflectância da fração terra fina seca ao ar (< 2 mm) [`reflectancia`]. Determinação em amostra em estado de umidade seco [`seco`]. Quantificada em ambiente controlado de laboratório [`lab`]. \{ESPECIFICAR A RESOLUÇÃO ESPECTRAL\} \{ESPECIFICAR ESPECTRORADIÔMETRO UTILIZADO\} \{INDICAR REFERÊNCIA DO MÉTODO\}\\
\hline
`reflectancia\_umido\_lab` & Reflectância da fração terra fina seca ao ar (< 2 mm) [`reflectancia`]. Determinação em amostra em estado de umidade seco [`seco`]. Quantificada em ambiente controlado de laboratório [`lab`]. \{ESPECIFICAR A RESOLUÇÃO ESPECTRAL\} \{ESPECIFICAR ESPECTRORADIÔMETRO UTILIZADO\} \{INDICAR REFERÊNCIA DO MÉTODO\}\\
\hline
`reflectancia\_umido\_estatico` & Reflectância do material do solo [`reflectancia`]. Determinação em amostra em estado de umidade seco [`seco`]. Quantificada em ambiente de campo utilizando equipemanto sensor estático [`estatico`]. \{ESPECIFICAR A RESOLUÇÃO ESPECTRAL\} \{ESPECIFICAR ESPECTRORADIÔMETRO UTILIZADO\} \{INDICAR REFERÊNCIA DO MÉTODO\}\\
\hline
`reflectancia\_umido\_movel` & Reflectância do material do solo [`reflectancia`]. Determinação em amostra em estado de umidade seco [`seco`]. Quantificada em ambiente de campo utilizando equipemanto sensor móvel [`movel`]. \{ESPECIFICAR A RESOLUÇÃO ESPECTRAL\} \{ESPECIFICAR ESPECTRORADIÔMETRO UTILIZADO\} \{INDICAR REFERÊNCIA DO MÉTODO\}\\
\hline
\end{tabular}

\hypertarget{harmonizauxe7uxe3o-1}{%
\section{Harmonização}\label{harmonizauxe7uxe3o-1}}

\hypertarget{refs}{}
\leavevmode\hypertarget{ref-AlmeidaEtAl2017}{}%
Almeida, Brivaldo Gomes de, João Herbert Moreira Viana, Wenceslau Geraldes Teixeira, and Guilherme Kangussu Donagemma. ``Densidade Do Solo.'' In \emph{Manual de Métodos de Análise de Solo}, edited by Paulo César Teixeira, Guilherme Kangussu Donagemma, Ademir Fontana, and Wenceslau Geraldes Teixeira, 65--75. Brasília: EMBRAPA, 2017.

\leavevmode\hypertarget{ref-Al-ShammaryEtAl2018}{}%
Al-Shammary, Ahmed Abed Gatea, Abbas Z. Kouzani, Akif Kaynak, Sui Yang Khoo, Michael Norton, and Will Gates. ``Soil Bulk Density Estimation Methods: A Review.'' \emph{Pedosphere} 28, no. 4 (August 2018): 581--96. \url{https://doi.org/10.1016/s1002-0160(18)60034-7}.

\leavevmode\hypertarget{ref-CasanovaEtAl2016}{}%
Casanova, Manuel, Elizabeth Tapia, Oscar Seguel, and Osvaldo Salazar. ``Direct Measurement and Prediction of Bulk Density on Alluvial Soils of Central Chile.'' \emph{Chilean Journal of Agricultural Research} 76, no. 1 (March 2016): 105--13. \url{https://doi.org/10.4067/s0718-58392016000100015}.

\leavevmode\hypertarget{ref-GrossmanEtAl2002}{}%
Grossman, R. B., and T. G. Reinsch. ``2.1 Bulk Density and Linear Extensibility.'' In \emph{SSSA Book Series}, edited by Jacob H. Dane and G. Clarke Topp, 201--28. Madison, WI, USA: Soil Science Society of America, 2002. \url{https://doi.org/10.2136/sssabookser5.4.c9}.

\leavevmode\hypertarget{ref-ICRAFEtAl2010}{}%
ICRAF, and ISRIC. ``A Globally Distributed Soil Spectral Library Visible Near Infrared Diffuse Reflectance Spectra.'' Nairobi, Kenya: World Agroforestry Centre (ICRAF), 2010.

\leavevmode\hypertarget{ref-MathieuEtAl1998}{}%
Mathieu, Clément, and Françoise Pieltain. \emph{Analyse Physique Des Sols: Méthodes Choisies}. Paris: Lavoisier, 1998.

\leavevmode\hypertarget{ref-OliveiraEtAl2018b}{}%
Oliveira, Virlei Álvaro de, José Frutuoso do Vale Júnior, Carlos Ernesto Gonçalves Reynaud Schaefer, José Francisco Lumbreras, Maurício Rizzato Coelho, Valdinar Ferreira Melo, Milton César Costa Campos, Guilherme Resende Corrêa, Sebastião Barreiros Calderano, and João Carlos Ker. ``Capítulo 2. Solos Da Xi Reunião Brasileira de Classificação E Correlação de Solos (Rcc de Roraima).'' In \emph{Guia de Campo Da Xi Reunião Brasileira de Classificação E Correlação de Solos: RCC de Roraima}, First., 12--29. Brasília, DF: EMBRAPA, 2018.

\leavevmode\hypertarget{ref-RibeiroEtAl2018}{}%
Ribeiro, Eloi, Niels H. Batjes, and A. J. M. van Oostrum. \emph{World Soil Information Service (Wosis) - Towards the Standardization and Harmonization of World Soil Data. Procedures Manual 2018}. Wageningen, The Netherlands: ISRIC World Soil Information, 2018. \url{https://doi.org/10.17027/isric-wdcsoils.20180001}.

\leavevmode\hypertarget{ref-ViscarraRosselEtAl2016}{}%
Viscarra Rossel, R. A., T. Behrens, E. Ben-Dor, D. J. Brown, J. A. M. Demattê, K. D. Shepherd, Z. Shi, et al. ``A Global Spectral Library to Characterize the World's Soil.'' \emph{Earth-Science Reviews} 155 (2016): 198--230. \url{https://doi.org/10.1016/j.earscirev.2016.01.012}.

\end{document}
